%! LaTeX Vorlage
\documentclass[12pt, fleqn, captions=nooneline, titlepage, footsepline, headsepline, toc=sectionentrywithdots, listof=entryprefix, bibliography=totoc, parskip=half-]{scrartcl}

\usepackage{silence} %unnötige Warnungen unterdrücken
\WarningFilter{latex}{You have requested}
\WarningFilter{scrlayer-scrpage}{\headheight to low}
\WarningFilter{scrlayer-scrpage}{\footheight to low}
\WarningFilter{scrlayer-scrpage}{Very small head height detected}
\WarningFilter{fvextra}{}
\WarningFilter{lineno}{}


\usepackage{tocbasic}
\usepackage[ngerman]{babel}
%\usepackage[backend=biber, style=authortitle, isbn=false]{biblatex}

%! Das Inhaltsverzeichnis wird an dieser Stelle formatiert.
\RedeclareSectionCommands[beforeskip=-.1\baselineskip, afterskip=.1\baselineskip, tocindent=0pt, tocnumwidth=45pt]{section,subsection,subsubsection}

%! Formatierung aller Verzeichnisse
\renewcaptionname{ngerman}{\refname}{Quellenverzeichnis}
%\setuptoc{toc}{totoc}
\setuptoc{lof}{totoc}
\setuptoc{lot}{totoc}
\renewcommand*\listoflofentryname{\bfseries\figurename}
\BeforeStartingTOC[lof]{\renewcommand*\autodot{\space\space\space\space}}
\addtokomafont{captionlabel}{\bfseries}
\renewcommand*\listoflotentryname{\bfseries\tablename}
\BeforeStartingTOC[lot]{\renewcommand*\autodot{\space\space\space\space}}

%? Anhangsverzeichnis
\providecaptionname{ngerman}{\listofatocentryname}{Anhang}
\DeclareTOCStyleEntry[level=0, indent=50pt, numwidth=5cm, entrynumberformat={Anlage}]{tocline}{atoc}

\makeatletter
\AfterTOCHead[atoc]{\let\if@dynlist\if@tocleft}
\newcommand*{\useappendixtocs}{
  \renewcommand*{\ext@toc}{atoc}
  \RedeclareSectionCommands[tocindent=0pt]{section, subsection, subsubsection}
  \RedeclareSectionCommands[tocnumwidth=85pt]{section, subsection, subsubsection}
  \addtokomafont{sectionentry}{\mdseries}
  \renewcommand*\listoflofentryname{\mdseries}
  \renewcommand{\thesection}{\arabic{section}}
  \renewcommand{\@dotsep}{10000}
  }
\newcommand*{\usestandardtocs}{
  \renewcommand*{\ext@toc}{toc}
  }
\makeatother

%! Ermöglicht die Ausgabe des aktuellen Titels
\usepackage{nameref}
\makeatletter
\newcommand*{\currentname}{\@currentlabelname}
\makeatother
\usepackage{xcolor}


\usepackage[normalem]{ulem}
\useunder{\uline}{\ul}{}
%! Latex-Packages, die verwendet werden
\usepackage{amsmath}
\usepackage{amssymb}
\usepackage{amsthm}
\usepackage{tabularx}
\usepackage{multirow}
\usepackage{setspace} 
\usepackage{booktabs}
\usepackage{svg}
\usepackage{graphicx}
\usepackage{float}
\usepackage[a4paper,lmargin={2.5cm},rmargin={2.5cm},tmargin={2cm},bmargin={2cm}]{geometry}
\usepackage{lineno}
\usepackage{csquotes}

%! Code Integration im Dokument
%? Inklusive Erzeugung eines Custom Enviroments für Programmcodes
\usepackage[outputdir=latex.out]{minted}
%\SetupFloatingEnvironment{listing}{name=Program code}
%\SetupFloatingEnvironment{listing}{listname=List of Program Code}


\DeclareNewTOC[
  type=code,                         % Name der Umgebung
  types=codes,                       % Erweiterung (\listofschemes)
  float,                               % soll gleiten
  floatpos=H,
  tocentryentrynumberformat=\bfseries,                        % voreingestellte Gleitparameter
  name=Code,                         % Name in Überschriften
  listname={Programmcodeverzeichnis}, % Listenname
  % counterwithin=chapter
]{loc}
\setuptoc{loc}{totoc}

\renewcommand*{\codeformat}{%
  \codename~\thecode%
  \autodot{\space\space\space}
}
\BeforeStartingTOC[loc]{\renewcommand*\autodot{\space\space\space\space}}


%! Lang ersehntes Formelumgebung / Formelverzeichnis
\DeclareNewTOC[
  type=formel,                         % Name der Umgebung
  types=formeln,                       % Erweiterung (\listofschemes)
  float,                               % soll gleiten
  tocentryentrynumberformat=\bfseries,                        % voreingestellte Gleitparameter
  name=Formel,                         % Name in Überschriften
  listname={Formelverzeichnis}, % Listenname
  % counterwithin=chapter
]{lom}
\setuptoc{lom}{totoc}
\renewcommand*{\formelformat}{%
  \formelname~\theformel%
  \autodot{\space\space\space}
}
\BeforeStartingTOC[lom]{\renewcommand*\autodot{\space\space\space\space}}


%! Schriftart
\usepackage{helvet}
\usepackage{microtype}
\renewcommand{\familydefault}{\sfdefault}
%! Hyperref und PDF-meta
\usepackage[hidelinks]{hyperref}
\usepackage[numbered]{bookmark}
\usepackage{acronym}
\usepackage{enumitem}

%Kein Umbruch auf der titelseite
\let\endtitlepage\relax

\renewcommand*{\aclabelfont}[1]{\acsfont{#1}} %Abkürzungsverzeichnis - Formatierung
 %zB zum korrekten anzeigen von BASH-Befehlen


%! Caption um Tabellen und Abbildungen richtig zu beschriften
\usepackage{caption}
% Captions linksbündig auch wenn einzeilig
\captionsetup{
  labelsep=none,
  justification=raggedright,
  singlelinecheck=false
}
\renewcommand*{\figureformat}{%
  \figurename~\thefigure%
  \autodot{\space\space\space}
}
\renewcommand*{\tableformat}{%
  \tablename~\thetable%
  \autodot{\space\space\space}
}


%! Formatierung der Fußnotenzitate / Literaturverzeichnis


%! Kopf- und Fußzeile
\usepackage[automark]{scrlayer-scrpage} 
\pagestyle{scrheadings} 
\clearmainofpairofpagestyles 
\clearplainofpairofpagestyles
\renewcommand*{\sectionmarkformat}{}
\rohead{\textnormal{\headmark}}
%\lohead{\includegraphics[height=8mm]{bilder/logo.png}}
\lohead{}
\lofoot{}
\cofoot{} 
\rofoot{\textnormal{Seite}~\pagemark}
\makeatletter
\usepackage{geometry}
\setlength{\footheight}{18.85002pt}
\geometry{a4paper,
          left=25mm,right=25mm,top=20mm,bottom=20mm,
          includehead=false,
          includefoot=false,
          headheight = \baselineskip,
          headsep = \dimexpr\Gm@tmargin-\headheight-10mm,
          footskip = \dimexpr\Gm@bmargin-10mm,
          %showframe,
          bindingoffset=0mm}
\setlength{\footheight}{\baselineskip}
\makeatother

%! Versuch von besseren Seitenumbrüchen
\clubpenalty = 10000
\widowpenalty = 10000
\displaywidowpenalty = 10000
\widowpenalties= 3 10000 10000 150

\linespread{1.3}
\newcommand\frontmatter{%
    \cleardoublepage
  %\@mainmatterfalse
  \pagenumbering{Roman}}

\newcommand\mainmatter{%
    \cleardoublepage
 % \@mainmattertrue
  \pagenumbering{arabic}}

\newcommand\backmatter{%
  \if@openright
    \cleardoublepage
  \else
    \clearpage
  \fi
 % \@mainmatterfalse
   }

%! INHALTSVERZEICHNIS / ANHANGSSVERZEICHNIS
\DeclareNewTOC[%
  %owner=\jobname,
  tocentrystyle=tocline,
  tocentryentrynumberformat=\PrefixBy{Anhang},
  listname={Anhangverzeichnis}% Titel des Verzeichnisses
]{atoc}% Dateierweiterung (a=appendix, toc=table of contents)


% ? Ab hier beginnen eigene Befehle um den Umgang zu erleichtern   
\newcommand{\logisch}[1]{$``#1``$}
\newcommand{\bild}[4][1.0]{\begin{figure}[H]
                      \centering
                      \includegraphics[width=#1\columnwidth]{bilder/#2}
                      \caption{#3}
                      \label{fig:#4}
                      \end{figure}}
\newcommand{\fullref}[1]{(siehe Kapitel~\ref{#1}~-~\nameref{#1})}
\newcommand{\literef}[1]{Kapitel~\ref{#1}~-~\nameref{#1}}
\newcommand{\aref}[1]{\emph{\hyperref[{#1}]{Anhang \ref{#1}}}}
\newcommand{\bref}[1]{\emph{\hyperref[{fig:#1}]{Abbildung \ref{fig:#1}}}}
\newcommand{\striche}[1]{\glqq #1\grqq{}}
\newcommand{\vglink}[2]{\footnote{\hspace{0.5em}vgl.~\href{#1}{#1}~(#2)}}
\newcommand{\python}[1]{\mintinline{python}{#1}}
\newcommand{\svg}[4][1.0]{\begin{figure}[H]
  \centering
  \includesvg[width=#1\columnwidth,inkscapelatex=false]{bilder/#2}
  \caption{#3}
  \label{fig:#4}
  \end{figure}}
\newcommand{\fn}[1]{\footnote{\hspace{0.5em}#1}}

\graphicspath{./bilder/}

%Für diese Version des Dokuments spezielle Anpassungen
\usepackage{array}
\newcolumntype{P}[1]{>{\centering\arraybackslash}p{#1}}
\definecolor{green}{RGB}{0, 153, 0}
\newcommand{\lat}[2]{(\textit{Lateinisch \glqq{}#1\grqq{} = #2})}
\newcommand{\merk}[1]{\glqq{}\textit{#1}\grqq{}}
\newcommand{\rot}[1]{\textcolor{red}{#1}}
\newcommand{\gruen}[1]{\textcolor{green}{#1}}
\newcommand{\schwarz}[1]{\textcolor{black}{#1}}
\newcommand{\menge}[1]{\{ #1 \}}
\newcommand{\komplement}[1]{\overline{#1}^E}
\tikzstyle{every picture}+=[remember picture]
\everymath{\displaystyle}



\lohead{\includesvg[height=8mm]{bilder/ba_glauchau_logo}}

\begin{document}
\frontmatter

\title{Vorlesungsskript Ingenieursmathematik 4DE19-1-IMA}

\author{Dominic Sczyrba}
\date{14. September 2021}

\maketitle
\tableofcontents

\mainmatter
\setlength{\mathindent}{0cm}

\section{Allgemeine Grundlagen Elementarer Mathematik}
    \subsection{Axiome, Definitionen und Sätze (vgl. ET-LV)}
        \label{sec:grundlagen}
        \subsubsection{Axiome}
            \begin{itemize}[leftmargin=*]
                \item Grundannahmen die aus der unmittelbaren Anschauung resultieren
                \item Einleichtentes Prinzip, das keines Beweises bedarf
                \item Formale Eigenschaft
                \item Eine unabgeleitete Aussage
            \end{itemize}
            \textit{Beispiele:}
            \begin{itemize}[leftmargin=*]
                \item Newtonsche Axiome
                \item Grundgesetz der Elektrostatik
                \item Kommutativgesetz \lat{commutare}{vertauschen}\\
                    \vspace{-0.8cm}
                    \begin{eqnarray*}
                        a \ +     \ b &=& b \ +     \ a \\
                        a \ \cdot \ b &=& b \ \cdot \ a \\
                        a \ \land \ b &=& b \ \land \ a \\
                        a \ \lor  \ b &=& b \ \lor  \ a
                    \end{eqnarray*}
                \item Assoziativgesetz  \lat{soziare}{zusammengehörig}
                \item Distributivgesetz \lat{distribuere}{verteilen}
                \item Geometrie - Euklid \\
                    \merk{Durch einen Punkt außerhalb einer Gerade gibt es nur eine Gerade, die sich mit der Gegebenen nicht schneidet.}
                    \\ (Zeichnung parallele Linien und Zeichnung gekrümmte Ebene einfügen)
            \end{itemize}
        \subsubsection{Definitionen}
            \begin{itemize}[leftmargin=*]
                \item (sinnvolle) Festlegungen
                \item ebenfalls nicht beweisbar
            \end{itemize}
            \textit{Beispiele:}
            \begin{itemize}[leftmargin=*]
                \item $\varrho = \frac{dQ}{dV}$
                \item $i = \frac{d}{dt} Q(t)$
                \item $R = \frac{U}{I}$
                \item   \begin{tabular}[t]{@{}l l}
                        Potenzieren     &  $a^b  \ \equiv \ \overbrace{a \cdot a \cdot a \cdot a \cdot \ldots}^{b \text{-Faktoren}} $ \\
                                        &  $2^{4}\ \equiv \ {2 \cdot 2 \cdot 2 \cdot 2 \cdot 2 \cdot = 32}$ \\ 
                                        &  $a^0 \ = \ 1$ \\
                        $\implies$ Idee: &  $a^b = \underbrace{a \cdot a \cdot a \cdot a}_{b \text{-Faktoren}} \cdot 1$ \\  
                        \end{tabular}
            \end{itemize}
        \subsubsection{Sätze}
            \begin{itemize}[leftmargin=*]
                \item wahrheitsfähige und beweispflichtige Aussagen
            \end{itemize}
            \textit{Beispiel: } Wenn eine Zahl durch $8$ teilbar ist, dann ist sie auch durch $4$ teilbar. 

            \begin{tabular}{@{}l l}
                Beweis: & Es sei $x\in\mathbb{N}$, wenn $\frac{x}{8} = y$ und $y\in\mathbb{N}$. \\
                & $\implies \frac{x}{2 \cdot 4} = y | \cdot 2$ dann ist $\frac{x}{4} = 2 \cdot y$ d.h. ebenfalls $\in\mathbb{N}$. \\
            \end{tabular}

            \rot{$\Rightarrow$ Merke:} \merk{Formeln gibt es weder in Mathematik noch in Physik.}
    \subsection{Mengenlehre}
        lt. Georg Cantor (1845 - 1918)
            \subsubsection{Begriffe und Darstellungen}
                \rot{Definition Menge:} \merk{Zusammenfassung von bestimmten wohlunterscheidbaren Objekten (Elementen) unserer Anschauung oder unseres Denkens}

                \textit{Beispiele:}
                    \begin{itemize}[leftmargin=*]
                        \item Menge aller Computer
                        \item Menge aller Programmiersprachen
                        \item Menge aller Lösungen einer Gleichung
                        \item Menge aller natürlichen Zahlen
                        \item Menge aller Gedanken
                    \end{itemize}

                \textit{Wichtig:}
                    \begin{itemize}[leftmargin=*]
                        \item mindestens \bfseries eine \mdseries gemeinsame Eigenschaft
                        \item für die Entscheidung, ob ein Element zu einer Menge gehört oder nicht, gibt es die Symbolik: \\
                            $x \in M \quad x \notin M $ \} ist stets eindeutig (deterministisch)
                        \item \rot{Achtung:} Die Elemente einer Menge müssen nicht geordnet sein, z.B. Die Menge aller Ebenenfiguren (Kreise, Quadrate, Rechtecke)
                        \item \textbf{Darstellungsformen:}
                        \item[] \begin{tabular}{@{}l p{9cm}}
                            a) aufzählende Form             & $M_1 = \{1, \ 2, \ 3\} $  \\
                                                            & $M_2 = \{\blacklozenge, \ \blacksquare, \ \blacktriangle \} $  \\
                                                            & $M_3 = \{\text{Max}, \ \text{Moritz}\} $ \\
                                                            & $M_4 = \{true, \ false\} $ \\
                            b) fest eingeführte Buchstaben  &  $ x \ \in \ \mathbb{N}$\\ 
                            c) beschreibende Form           & $A = \{a \underbrace{|}_{mit} a \text{~besitzt die Eigenschaft~} E\}$\\
                                                            & $B = \{b | b \text{~ist ein weiblicher Vorname} \} $ \\
                                                            & $C = \{c | c \text{~ist eine gerade Zahl} \} $ \\
                                                            & $D = \{d \in \mathbb{N} | d^2 - 3d + 2 = 0 \} $\\
                            d) Venn - Diagramm              & Skizze vom Venn Diagramm einfügen \\
                            e) Zahlenstrahl                 & Skizze vom Zahlenstrahl \\
                            f) geometrische Gebilde         & Ein Kreis ist die Menge aller Punkte, die in der Ebene von einem Mittelpunkt $M$ den gleichen Abstand $r$ besitzen.\\
                            g) als Tabelle                  & DE19 \quad $\underbrace{\text{Matrikelnummer \ Name \ Vorname}}_{\text{eine Relation}}$ \\
                            \end{tabular}
                        \item \textbf{Kennzeichnung der leeren Menge}
                        \item[] $M \ = \ \{\}$ oder $M \ = \ \varnothing$
                        \item \textbf{Kennzeichnung von endlicher / unendlicher Menge}
                        \item[] $ M \ = \ \{ 1,\ 2,\ 3, \ 4, \ \ldots, \ 20\} $ \\
                                $ M \ = \ \{ 1,\ 2,\ 3, \ 4, \ \ldots \} $
                        \item \textbf{Kennzeichnung der Einheitsmenge $\mathbb{E}$}
                        \item[] Venn-Diagramm mit Einheitsmenge einfügen
                        \item \textbf{Kennzeichnung der Anzahl der Elemente}
                        \item[] $|m| $ \\ $|\{\rhd,\ \triangleleft,\ \triangleright \}| \ = \ 3$
                    \end{itemize}
                \subsubsection{Beziehungen zwischen Mengen}
                \begin{itemize}[leftmargin=*]
                    \item[a)] $A = B$ \quad (auch $A\equiv B$) z.B. $\{1, 2\} = \{2, 1\}$ 
                    \item[b)] $\underbrace{A\subseteq B}_{Teilmenge}$ bzw. $\underbrace{A \subset B}_{echte~Teilmenge}$
                    \item[] z.B. $\{1,\ 2\} \subset \{1,\ 5,\ 2\}$
                    \item[] \begin{tikzpicture}
                        \filldraw[color=magenta, fill=magenta!10, thick] (0,0) ellipse (2.5 and 1.5);
                        \filldraw[color=cyan, fill=cyan!10, thick] (-0.5, 0) ellipse (1.5 and 1.0);
                        \coordinate[label=center:$\textcolor{cyan}{2}$] (M) at (0, -0.5);
                        \coordinate[label=center:$\textcolor{cyan}{1}$] (M) at (-1, 0.5);
                        \coordinate[label=center:$\textcolor{magenta}{5}$] (M) at (1.5, 0);
                        \end{tikzpicture} 
                    \item[] \underline{Satz:} Wenn $A \subseteq B$ und $B \subseteq A \Leftrightarrow A = B$
                    \item[c)] Komplementmenge (Bezugsmenge oder Grundmenge)
                    \item[] \begin{tikzpicture}
                        \filldraw[color=red,fill=red!10, thick] (0, 0) rectangle (4, 3);
                        \draw[color=blue, fill=blue!5, thick] (2, 1.5) circle (1.2);
                        \coordinate[label=center:$\textcolor{blue}{M}$] (M) at (2, 1.5);
                        \coordinate[label=center:$\textcolor{red}{E}$] (E) at (0.3, 0.3);
                        \end{tikzpicture}\quad \textcolor{red}{$\overline{\textcolor{blue}{M}}^E$} Komplement zu \textcolor{blue}{M}. bzgl. \textcolor{red}{E}
                    \item[d)] Potenzmenge  
                    \item[] Potenzmenge $P(M)$ $\Rightarrow$ Menge aller Teilmengen von $M$
                    \item[] \textit{Beispiel:}\hphantom{sss} $M = \{1, \ 2, \ 3\}$
                    \item[] \hphantom{Beispiel}\hphantom{sss} $M = \Big\{\{1\}, \{2\}, \{3\}, \{1, 2\}, \{1, 3\}, \{2, 3\}, \{1, 2, 3\}\Big\} + \varnothing $
                    \item[] \rot{Satz:}\hphantom{sss} Wenn $n \ = \ |m|$, dann $|P(m)| \ = \ 2^n$
                \end{itemize}
                \newpage
            \subsubsection{Mengenoperationen}
                \begin{itemize}[leftmargin=*]
                    \item[1.] \textbf{Schnittmengenbildung = Durchschnitt}
                    \item[] $A \ \textcolor{red}{\cap} \ B \ = \ C \implies C = \{c|c \in A \ \text{\underline{und}} \ c \in B\}$ \\
                            $A~geschniten~B$
                    \item[] \textit{Venn-Diagramm}
                    \item[] \begin{minipage}{0.6\columnwidth}\begin{tikzpicture}
                        \draw[color=blue, thick, rotate=-15] (0,0) ellipse (2.5 and 1.5);
                        \draw[color=green, thick, rotate=5] (2.5,-0.5) ellipse (2.0 and 1.5);
                        \draw[color=red, thick] (-1.5, -2) ellipse (1.0 and 0.5);
                        \node (1) at (-0.5, 0.9) {1};
                        \draw[color=magenta, thick] (1) circle (0.3) node[anchor=west] {\textcolor{magenta}{~~C}};
                        \node (2) at (1.4, 0.3) {2};
                        \node (3) at (1.3, -0.6) {3};
                        \node (4) at (3.4, 0) {4};
                        \node (5) at (-1.0, -2) {5};
                        \node (6) at (-2.0, -2) {6};
                        \node (A) at (0, 1.8) {\textcolor{blue}{A}};
                        \node (B) at (4.8, 0) {\textcolor{green}{B}};
                        \node (C) at (-2, -2.7) {\textcolor{red}{D}};
                        \end{tikzpicture}\end{minipage}
                        \begin{minipage}{0.4\columnwidth}
                            $A \cap B = \{2, 3\}$ \\
                            $A \cap C = \{1\}$ \\
                            $A \cap D = \varnothing$ \\
                            $B \cap C = \varnothing$ \\
                            $B \cap D = \varnothing$ \\
                            $C \cap D = \varnothing$ \\
                        \end{minipage}
                    \item[] \textit{Axiome}
                        \begin{align*}
                            A \cap B &= B \cap A & \text{Kommutativität} \\
                            A \cap B \cap C &= (A \cap B) \cap C & \text{Assoziativität}\\
                                            &= A \cap (B \cap C) & \text{Assoziativität}\\
                            A \cap A &= A &    \text{Idempotenz} \\
                        A \cap \varnothing &= \varnothing & \text{Extremalgesetz} \\
                        A \cap E &= A & \text{Neutralitätsgesetz} \\        
                        \end{align*}
                    \item[] \textit{Verallgemeinerte Symbolik}
                        \begin{align*}
                        \textcolor{red}{\bigcap_{\textcolor{green}{i=1}}^{\textcolor{green}{n}}} M_{\textcolor{green}{i}} &= M_1 \cap M_2 \cap M_3 \cap \ldots \cap M_n \\
                        \textcolor{red}{\bigcap_{\textcolor{green}{i=1}}^{\textcolor{green}{\infty}}} M_{\textcolor{green}{i}} &= \lim_{n \to \infty} \textcolor{red}{\bigcap_{\textcolor{green}{i=1}}^{\textcolor{green}{n}}} M_{\textcolor{green}{i}} = M_1 \cap M_2 \cap M_3 \cap \ldots \cap M_n \\        
                        \end{align*}
                    \item[2.] \textbf{Vereinigungsmenge}
                    \item[] $A \ \textcolor{red}{\cup} \ B = C \Longleftrightarrow C = \{c|c \in A ~\text{\underline{oder}}~ c \in B\}$
                    \item[] \textit{Axiome} \begin{align*}
                        A \cup B &= B \cup A & \text{Kommutativität} \\
                        A \cup B \cup C &= (A \cup B) \cup C & \text{Assoziativität} \\
                                        &= A \cup (B \cup C) & \text{Assoziativität} \\
                        A \cup A &= A & \text{Itempotenz} \\
                        A \cup \varnothing &= A & \text{Extremalgesetz} \\
                        A \cup E &= E & \text{Neutralitätsgesetz} \\
                        \end{align*}
                    \item[3.] \textbf{Differenzmenge}
                    \item[] $A \ \textcolor{red}{\setminus} \ B = C \Longleftrightarrow C = \{c|c \in A ~\text{und}~ c \notin B \}$
                    \item[] $A \ \setminus \ B = A \cap B^E$
                    \item[] Differenzmenge vervollständigen!   
                    \item[4.] \textbf{Komplementbildung}
                    \item[] $\textcolor{red}{\overline{\textcolor{black}{A}}^E} \ = \ C \Longleftrightarrow C = \{c|c \notin A \}$
                    \item[] umgeformt: $\overline{A}^E \ = \ E \ \setminus \ A$
                    \item[] \textit{Axiome:} \begin{eqnarray*}
                            \overline{\overline{A}^E}^E \ &=& \ A \\
                            \overline{E}^E \ &=& \ \varnothing \\
                            \overline{\varnothing}^E \ &=& \ E \\
                        \end{eqnarray*}
                    \item[5.] \textbf{Symmetrische Differenz}
                    \item[] $A \ \rot{\triangle} \ B \ = \ C \Longleftrightarrow C = \{c|c \in (A \setminus B) ~\text{\underline{oder}}~ c \in (B     \setminus A)\}$
                    \item[] Venn-Diagramm einfügen
                    \item[6.]\textbf{Komplement der symmetrischen Differenz}
                    \item[] $A \ \rot{\overline{\triangle}} \ B$
                    \item[] Venn-Diagramm einfügen 
                \end{itemize}
            \subsubsection{Weitere Rechenregeln (Axiome)} 
                \begin{itemize}[leftmargin=*]
                    \item \textbf{Distributivgesetz} \\
                    \begin{minipage}{0.5\columnwidth}
                        ~ \\
                        Der Vereinigung über  dem Durchschnitt
                        \begin{equation*}
                            \tikz[baseline]{\node[anchor=base] (t1) {$ A$};} \cup (
                            \tikz[baseline]{\node[anchor=base] (t2) {$B$}; } \cap
                            \tikz[baseline]{\node[anchor=base] (t3) {$C$}; } ) = (A\cup B) \cap (A\cup C)
                        \end{equation*}
                        \begin{tikzpicture}[overlay]
                            \path[->] (t1) edge [color=red, out=-60, in= 235] (t2);
                            \path[->] (t1) edge [color=red, out=-75, in=235] (t3);
                        \end{tikzpicture}
                    \end{minipage}
                    \begin{minipage}{0.5\columnwidth}
                        ~ \\
                        Des Durchschnitts über der Vereinigung
                        \begin{equation*}
                            \tikz[baseline]{\node[anchor=base] (t1) {$ A$};} \cap (
                            \tikz[baseline]{\node[anchor=base] (t2) {$B$}; } \cup
                            \tikz[baseline]{\node[anchor=base] (t3) {$C$}; } ) = (A\cap B) \cup (A\cap C)
                        \end{equation*}
                        \begin{tikzpicture}[overlay]
                            \path[->] (t1) edge [color=red, out=-60, in= 235] (t2);
                            \path[->] (t1) edge [color=red, out=-75, in=235] (t3);
                        \end{tikzpicture}
                    \end{minipage}
                    \item[] \textit{Beispiele:}
                    \item[] 
                    \begin{minipage}{0.5\columnwidth}
                    \begin{tikzpicture}
                        \draw[color=green, thick, rotate=-20] (0,0) ellipse (2.0 and 1.5);
                        \node (A) at +(0, 1.8) {\textcolor{green}{A}};
                        \draw[color=red, thick, rotate=5] (2.5,0) ellipse (2.0 and 1.5);
                        \node (B) at (4.8, 0) {\textcolor{red}{B}};
                        \draw[color=blue, thick](1.25, -1.75) ellipse (2 and 1.5);
                        \node (C) at (1.25, -3.5) {\textcolor{blue}{C}};
                        \node (1) at (-0.9, 0.5) {1};
                        \node (2) at (1.0, 0.4){2};
                        \node (3) at (3, 0.6){3};
                        \node (4) at (0.4, -0.9){4};
                        \node (5) at (1.3, -0.6){5};
                        \node (6) at (2.3, -0.8){6};
                        \node (7) at (1.3, -2.2){7};
                        \end{tikzpicture}
                    \end{minipage}
                    \begin{minipage}{0.5\columnwidth}
                        $A \cup (B \cap C)$ \\
                        $B \cap C = \{ 5, 6\}$ \\
                        $\implies \{1, 2, 4, 5\} \cup \{5, 6\} = \{1, 2, 4, 5, 6\}$ \\
                        ~ \\
                        $(A \cup B) \cap (A \cup C) $\\
                        $ A \cup B = \menge{1, 2, 3, 4, 5, 6}$ \\
                        $ A \cup C = \menge{1, 2, 4, 5, 6, 7}$ \\
                        $\implies (A \cup B) \cap (A \cup C) = \menge{1, 2, 4, 5, 6}$ \\
                        ~ \\
                        $A \cap (B \cup C)$ \\
                        $B \cup C = \menge{2, 3, 4, 5, 6, 7}$ \\
                        $A = \menge{1, 2, 4, 5}$ \\
                        $\implies A \cap (B \cup C) = \menge{2, 4, 5}$ \\
                        ~ \\
                        $(A \cap B) \cup (A \cap C)$ \\
                        $A \cap B = \menge{2, 5}$ \\
                        $A \cap C = \menge{4, 5}$ \\
                        $\implies (A \cap B) \cup (A \cap C) = \menge{2, 4, 5}$
                    \end{minipage}
                \item \textbf{Komplementärgesetze} 
                \item[] $A \cap \komplement{A} = \varnothing$ 
                \item[] $A \cup \komplement{A} = \mathbb{E}$ 
                \item \textbf{Absorptionsgesetze}
                \item[] $A \cup (A \cap B) = A $
                \item[] $A \cup (\komplement{A} \cap B) = A \cup B $
                \item[] $A \cap (A \cup B) = A $
                \item[] $A \cap (\komplement{A} \cup B) = A \cap B $
                \item \textbf{Gesetze von De Morgan} 
                \item[] $\komplement{A \cup B} = \komplement{A} \cap \komplement{B}$
                \item[] $\komplement{A \cap B} = \komplement{A} \cup \komplement{B}$
                \item[] $ \left.
                    \begin{array}{c}  \komplement{\bigcup_{i=1}^n M_i} = \bigcap_{i=1}^n \komplement{M_i} \\
                                \komplement{\bigcap_{i=1}^n M_i} = \bigcup_{i=1}^n \komplement{M_i} \\
                    \end{array}%
                     \right\} \text{Dualitätsprinzip} $ %
                \item[] $A \setminus (B \cup C) = (A \setminus B) \cap (A \setminus C)$
                \item[] $A \setminus (C \cap C) = (A \setminus B) \cup (A \setminus C)$
                \end{itemize}
            \subsubsection{Ausgewählte Beispiele}
                \begin{itemize}[leftmargin=*]
                    \item[1.] $geg.\quad 4~$ Mengen $A \quad B \quad C \quad D$
                    \item[]   $ges.\quad A \cup (B \setminus A)$
                        \begin{align*}
                            A \cup (B \setminus A) & =A \cup (B \cap \komplement{A}) & \rightarrow X \setminus Y = X \cap \komplement{Y} \\
                                                   & =(A \cup B) \cap (A \cup \komplement{A})& \rightarrow \text{~Distributivgesetz}\\
                                                   & =(A \cup B) \cap E & \downarrow X \cap E = X \\
                            A \cup (B \setminus A) & =\underline{A \cup B} &\\
                        \end{align*}
                    \item[2.]  $ges.\quad (A \setminus B) \cap (B \setminus A)$
                        \begin{align*}
                                (A \setminus B) \cap (B \setminus A) & = (A \cap \komplement{B}) \cap (B \cap \komplement{A}) & \downarrow \text{~Assoziativgesetz} \\
                                & = A \cap \komplement{B} \cap B \cap \komplement{A} & \downarrow \text{~Umsortieren} \\
                                & = A \cap \komplement{A} \cap B \cap \komplement{B} & \\
                                & = \varnothing \cap \varnothing & \\
                                (A \setminus B) \cap (B \setminus A) & = \underline{\varnothing} & \\
                        \end{align*}
                    \item[3.] $ges.\quad (A \cup B) \setminus B$ %
                        \begin{align*}
                            (A \cup B) \setminus B & = (A \cup B) \cap \komplement{B} \\
                                                   & = (A \cap \komplement{B}) \cup (B \cap \komplement{B}) \\
                                                   & = (A \cap \komplement{B}) \cup \varnothing \\
                            (A \cup B) \setminus B & = \underline{A \cap \komplement{B}} \\                      
                        \end{align*}
                    \item[4.] $ges.\quad (A \setminus B) \cup (A \setminus \komplement{B})$ %
                        \begin{align*}
                            (A \setminus B) \cup (A \setminus \komplement{B}) & = (A \cap \komplement{B}) \cup (A \cap B) \\
                            & = (A \cup A) \cap (A \cup B) \cap (\komplement{B} \cup A) \cap (\komplement{B} \cup B) \\
                            & = \underbrace{A \cap (A \cup B)}_{Absorptionsgesetz} \cap (\komplement{B} \cup A) \underbrace{\cap \ E}_{entfaellt} \\
                            & = \underbrace{A \cap (\komplement{B} \cup A)}_{Absorptionsgesetz} \\
                            (A \setminus B) \cup (A \setminus \komplement{B}) & = \underline{A} \\
                        \end{align*}
                    \item[5.] $ges.\quad (A \cap B) \cup (A \cap \komplement{B})$  
                        \begin{align*}
                            (A \cap B) \cup (A \cap \komplement{B}) & = (A \cup A) \cap (A \cup \komplement{B}) \cap (B \cup A) \cap (B \cup \komplement{B}) \\
                                                                    & = A \cap (A \cup \komplement{B}) \cap (B \cup A) \cap E \\
                                                                    & = A \cap (A \cup B) \\
                            (A \cap B) \cup (A \cap \komplement{B}) & = \underline{A}                                       
                        \end{align*}
                    \item[6.] $ges.\quad (A \cap B) \cup (A \cap \komplement{B})$  
                        \begin{align*}
                            (A \cap B) \cup (A \cap \komplement{B}) & = A \cap (B \cup \komplement{B}) \\
                        \end{align*}
                    \item[7.] $ges.\quad (A \cap \komplement{D}) \cup (A \cap B \cap D) \cup (A \cap D)$  
                        \begin{align*}
                            (A \cap \komplement{D}) \cup (A \cap B \cap D) \cup (A \cap D) & = \Big[ A \cap (\komplement{D} \cup D)\Big] \cup (A \cap B \cap D)\\
                            & = [A \cap E] \cup (A \cap B \cap D) \\
                            & = A \cup (A \cap B \cap D) \\
                            & = (A \cup A) \cap (A \cup B) \cap (A \cup D) \\
                            & = A \cap (A \cup B) \cap (A \cup D) \\
                            & = A \cap (A \cup D) \\
                            (A \cap \komplement{D}) \cup (A \cap B \cap D) \cup (A \cap D) & = \underline{A} \\
                        \end{align*}
                    \item[8.] $ges.\quad (A \cap B) \cup (A \cap \komplement{B}) \cup (\komplement{A} \cap \komplement{B})$  
                        \begin{align*}
                            (A \cap B) \cup (A \cap \komplement{B}) \cup (\komplement{A} \cap \komplement{B}) & = A \cap (B \cup \komplement{B}) \cup (\komplement{A} \cap \komplement{B}) \\
                            & = A \cup (\komplement{A} \cap \komplement{B}) \\
                            & = (A \cup \komplement{A}) \cap (A \cup \komplement{B}) \\
                            & = E \cap (A \cup \komplement{B}) \\
                            (A \cap B) \cup (A \cap \komplement{B}) \cup (\komplement{A} \cap \komplement{B}) & = \underline{A \cup \komplement{B}} \\
                        \end{align*}
                    \item[9.] $ges.\quad \komplement{A \triangle B}$  
                        \begin{align*}
                            \komplement{A \triangle B} & = \komplement{(A \setminus B) \cup (B \setminus A)} & \downarrow \text{~de Morgan}\\
                            & = \komplement{(A \setminus B)} \cap \komplement{(B \setminus A)} \\
                            & = \komplement{(A \cap \komplement{B})} \cap \komplement{(B \cap \komplement{A})} & \downarrow \text{~de Morgan} \\
                            & = (\komplement{A} \cup B) \cap (\komplement{B} \cup A) \\
                            & = \komplement{A} \cap \komplement{B} \cup \varnothing \cup \varnothing \cup (B \cap A) \\
                            & = (\komplement{A} \cup B) \cap (\komplement{B} \cup A) \\
                        \end{align*}
                    \item[10.] Bei Gelegenheit weitere Beispiele übernehmen
                \end{itemize}
            \subsubsection{Abbildungen}
                \begin{itemize}[leftmargin=*]
                    \item \textbf{Definition:} Es seien $A, B$~Mengen, dann ist eine \gruen{Abbildung} die \rot{Menge der geordneten Paare $(a, b)$}~mit $\rot{a} \in A$~und $\rot{b} \in B$, die über eine \gruen{Abbildungsvorschrift} miteinander verknüpft sind.
                    \item[]\textit{Beispiel} \\
                        \begin{tikzpicture}
                            \node (M1) at (-2, 2.3) {\Large$\textcolor{blue}{A}$};
                            \node (M2) at (4, 2.3) {\Large$\textcolor{blue}{B}$};
                            \draw[color=blue, thick](-2, 0) circle (2.0);
                            \node (A) at (-2.8, 0.8) {\Huge$\textcolor{blue}{\triangle}$};
                            \node (B) at (-1.2, 0.8) {\Huge$\textcolor{blue}{\bigcirc}$};
                            \node (C) at (-2.8, -0.8) {\Huge$\textcolor{blue}{\square}$};
                            \node (D) at (-1.2, -0.8) {\Huge$\textcolor{blue}{\boxplus}$};
                            \draw[color=blue, thick](4, 0) circle (2.0);
                            \node (E) at (3.2, 1.0) {\Large$\textcolor{blue}{Max}$};
                            \node (F) at (3.6, 0.0) {\Large$\textcolor{blue}{Fritz}$};
                            \node (G) at (4, -1.0) {\Large$\textcolor{blue}{Franz}$};
                            \path[->] (B) edge [color=green, thick] (E);
                            \path[->] (B) edge [color=green, thick] (F);
                            \path[->] (D) edge [color=green, thick] (F);
                        \end{tikzpicture}
                    \item[] \gruen{Abbildung} $= \Big\{ \rot{(} \bigcirc ; Max\rot{)}; \rot{(} \bigcirc ; Fritz\rot{)} ; \rot{(} \boxplus ; Fritz\rot{)} \Big\}$ 
                    \item[] Die Menge aller \rot{a}, die ein \gruen{Abbild} \rot{a} besitzen, heißen Definitionsbereich der \gruen{Abbildung}.
                    \item[] \textit{Beispiel}
                    \item[] Definitionsbereich $= \{\square; \bigcirc\}$
                    \item[] Die Menge aller \rot{b}, die \gruen{Abbilder} von Elementen aus \rot{a} darstellen, heißt \rot{Wertebereich} der \gruen{Abbildung}, \textit{z.B.} Wertebereich $= \{Max; Fritz\}$.
                    \item[] \textit{Sprechweisen}
                        \begin{itemize}
                            \item[] \glqq \ $~\ldots$ eine Abbildung \gruen{aus A nach B} \grqq 
                            \item[] \glqq \ $~\ldots$ eine Abbildung \gruen{aus A in B} \grqq 
                            \item[] \glqq \ $~\ldots$ eine Abbildung \gruen{von A auf B} \grqq  
                        \end{itemize}
                    \item[] \textit{Darstellung von Abbildungsvorschriften}
                        \begin{enumerate}
                            \item als Venn-Diagramm
                            \item als Tabelle
                            \item[] \textit{Beispiel} \\
                                \begin{tabular}{P{3cm}|P{3cm}}
                                    $A$             & $B$           \\\hline
                                    $\triangle$     & $\setminus$   \\
                                    $\bigcirc$      & $Max$         \\
                                    $\bigcirc$      & $Fritz$       \\
                                    $\square$       & $\setminus$   \\
                                    $\boxplus$      & $Fritz$       \\
                                    $\setminus$     & $Franz$       \\
                                \end{tabular}
                            \item als mathematisch-analytische Gleichung 
                            \item[] \textit{Beispiel} \\
                                \begin{minipage}{0.5\columnwidth}
                                    \begin{tabular}{P{3cm}|P{3cm}}
                                        $A$         & $B$       \\\hline
                                        $-3$        & $5$       \\
                                        $-2$        & $3$       \\ 
                                        $-1$        & $1$       \\
                                        $0$         & $-1$      \\
                                        $1$         & $1$       \\
                                        $2$         & $3$       \\
                                        $3$         & $5$       \\
                                    \end{tabular}
                                \end{minipage}
                                \begin{minipage}{0.5\columnwidth}
                                    \begin{equation*}
                                        b = \begin{cases}
                                            -2 \cdot a -1 & \text{für } a \leq 0 \\
                                            +2 \cdot a -1 & \text{für } a > 0 \\
                                        \end{cases}
                                    \end{equation*}
                                \end{minipage}
                            \item[] \textit{\rot{Aufgaben: } Finden Sie Abbildungsvorschriften.} \\
                                \begin{minipage}[t][8cm][t]{0.5\columnwidth}
                                    \begin{tabular}{P{3cm}|P{3cm}}
                                        $A$         & $B$       \\\hline
                                        $-3$ & $8$ \\
                                        $-2$ & $4$ \\
                                        $-1$ & $2$ \\
                                        $0$ & $1$ \\
                                        $1$ & $2$ \\
                                        $2$ & $4$ \\
                                        $3$ & $8$ \\
                                    \end{tabular}
                                    \begin{equation*}
                                        b = \begin{cases}
                                            2^{-a} & \text{für } a < 0 \\
                                            2^a    & \text{für } a \geq 0 \\ 
                                        \end{cases} = 2^{|a|}
                                    \end{equation*}
                                \end{minipage}
                                \begin{minipage}[t][8cm][t]{0.5\columnwidth}
                                    \begin{tabular}{P{3cm}|P{3cm}}
                                        $A$         & $B$       \\\hline
                                        $-3$ & $10$ \\
                                        $-2$ & $5$ \\
                                        $-1$ & $2$ \\
                                        $0$ & $1$ \\
                                        $1$ & $2$ \\
                                        $2$ & $5$ \\
                                        $3$ & $10$ \\
                                    \end{tabular}
                                    \begin{equation*}
                                        b = a^2 + 1  
                                    \end{equation*}
                                \end{minipage}                     
                        \end{enumerate}
                    \item[] \textit{Umkehrabbildung}
                \end{itemize}                    
\end{document}